% 第3章 3.1.1：l^2 上 A(x_n)=(x_n/n) 压缩单位球。
% (a) 在 span{e_1,e_2} 截面上，单位圆被映为半轴 1 与 1/2 的椭圆；
%     圆上取角度 0,40,...,320 度的九个点，其像为 (cos t, sin t / 2)。
% (b) 像集 A(B) 的第 n 个坐标满足 |y_n|<=1/n，条高为 1/n，n=1,...,16。
\begin{tikzpicture}[font=\figurefont\footnotesize]
\begin{scope}[scale=1.25]
  \draw[FigureGrid,line width=.4pt,-{Stealth[length=1.6mm]}] (-1.3,0) -- (1.35,0) node[right,text=FigureInk] {$e_1$};
  \draw[FigureGrid,line width=.4pt,-{Stealth[length=1.6mm]}] (0,-1.25) -- (0,1.3) node[above,text=FigureInk] {$e_2$};
  \draw[FigureBlue,line width=.9pt,dashed] (0,0) circle (1);
  \fill[FigureBlueFill] (0,0) ellipse (1 and .5);
  \draw[FigureBlue,line width=1.1pt] (0,0) ellipse (1 and .5);
  \foreach \t in {40,80,120,160,200,240,280,320} {
    \draw[FigureMuted,line width=.4pt,-{Stealth[length=1.2mm]}] ({cos(\t)},{sin(\t)}) -- ({cos(\t)},{0.5*sin(\t)+0.06*sign(sin(\t))});
    \fill[FigureBlue] ({cos(\t)},{sin(\t)}) circle (.9pt);
    \fill[FigureStroke] ({cos(\t)},{0.5*sin(\t)}) circle (.9pt);
  }
  \node[FigureBlue,anchor=south west] at (0.62,0.82) {$B$};
  \node[FigureStroke,fill=FigureBlueFill,inner sep=1pt] at (-0.42,0.2) {$A(B)$};
  \node[anchor=north,font=\figurefont\footnotesize] at (0,-1.38) {(a) 单位圆与它的像};
\end{scope}
\begin{axis}[at={(37mm,-15mm)},anchor=south west,width=52mm,height=36mm,
  ybar,bar width=1.5mm,xmin=0.3,xmax=16.7,ymin=0,ymax=1.1,
  axis x line=bottom,axis y line=left,
  xtick={1,4,8,12,16},ytick={0,0.5,1},
  xlabel={坐标序号$n$},ylabel={半轴长},
  xlabel style={font=\figurefont\scriptsize},ylabel style={font=\figurefont\scriptsize},
  tick label style={font=\figurefont\scriptsize,text=FigureInk},
  ymajorgrids,grid style={draw=FigureGrid,line width=.4pt},
  title={(b) 各坐标方向的半轴长$1/n$},title style={font=\figurefont\footnotesize}]
\addplot[fill=FigureBlue,draw=FigureBlue] coordinates {(1,1) (2,0.5) (3,0.3333) (4,0.25) (5,0.2) (6,0.1667) (7,0.1429) (8,0.125) (9,0.1111) (10,0.1) (11,0.0909) (12,0.0833) (13,0.0769) (14,0.0714) (15,0.0667) (16,0.0625)};
\end{axis}
\end{tikzpicture}
